\documentclass[12pt]{article}
\usepackage[utf8]{inputenc}
\usepackage[T1]{fontenc}
\usepackage{hyperref}
\title{Tala Language Resources for K-12 Curriculum Development}
\author{Compiled by Claude Sonnet 5 (claude-sonnet-5) for Andr\'as Kornai}
\date{}
\begin{document}
\maketitle
\section{Introduction}
Tala (ISO 639-3 tak; Glottocode tala1295) is a West Chadic language of Nigeria, in the Guruntum group of South Bauchi. Sources differ on its speakers: about 1,000 (1993) is the figure most often cited, and Joshua Project gives 2,200. Wikipedia describes it as unwritten, while Joshua Project lists a writing system. It should not be confused with Tal (tal), a separate language of Plateau State. Very little exists for Tala, and no educational materials were found.
\section{Primary Educational Resources}
None found.
\section{Linguistic and Reference Materials}
\begin{itemize}
\item Wikipedia overview of the Tala language: \url{https://en.wikipedia.org/wiki/Tala_language}
\item South Bauchi languages (classification): \url{https://en.wikipedia.org/wiki/South_Bauchi_languages}
\item Glottolog entry: \url{https://glottolog.org/resource/languoid/id/tala1295.bigmap.html}
\end{itemize}
\section{Religious and Community Resources}
\begin{itemize}
\item Joshua Project people-group profile (Tala in Nigeria): \url{https://joshuaproject.net/people_groups/15211}
\end{itemize}
\section{Digital and Audio Resources}
\begin{itemize}
\item Endangered Languages Project entry (rated threatened): \url{https://endangeredlanguages.com/lang/1021}
\end{itemize}
\section{Curriculum Development Context}
The language is not known to be taught in schools and is spoken mainly by older adults. Any curriculum would first need documentation: a wordlist, a practical orthography and community consultation.
\end{document}